\documentclass[12pt]{article}
\usepackage[utf8]{inputenc}
\usepackage[T1]{fontenc}
\usepackage{hyperref}
\title{Takelma Language Resources for K-12 Curriculum Development}
\author{Compiled by Claude Sonnet 5 (claude-sonnet-5) for Andr\'as Kornai}
\date{}
\begin{document}
\maketitle

\section{Introduction}
Takelma is a language of southwestern Oregon. It has no living first-language speakers; the last fluent speaker, Frances Johnson, died in 1934. The documentation rests mainly on Edward Sapir's 1906 fieldwork. The searches found a community revitalization effort, but no ready-made school curriculum.

\section{Primary Educational Resources}
The Cow Creek Band runs a Takelma language program with classes and books, according to a PBS item. No school curriculum or textbook was confirmed:
\url{https://www.pbs.org/video/cow-creek-tribal-takelma-language-i3enal}

\section{Linguistic and Reference Materials}
\begin{itemize}
\item Sapir's Takelma Texts (1909) and his Takelma grammar (1912) are the core sources, as described in the Wikipedia article: \url{https://en.wikipedia.org/wiki/Takelma_language}
\item Sapir's Takelma field notebooks (texts with translations, medicine formulas, grammar notes), ACLS collection, American Philosophical Society: \url{https://indigenousguide.amphilsoc.org/entry/40056}
\item University of Oregon Cascade article on descendants' revival work and a 2015 senior thesis on Takelma verb structure: \url{https://cascade.uoregon.edu/fall2015/humanities/a-language-within-a-language/index.html}
\end{itemize}

\section{Religious and Community Resources}
\begin{itemize}
\item Cow Creek Band Takelma language program (PBS, Southern Oregon Responds): \url{https://www.pbs.org/video/cow-creek-tribal-takelma-language-i3enal}
\end{itemize}

\section{Digital and Audio Resources}
Wikipedia mentions an English-Takelma dictionary being compiled from printed sources. No publisher or release date was found, and no audio resources were found in the searches.

\section{Curriculum Development Context}
A heritage curriculum would have to be built from Sapir's texts and grammar, with the Cow Creek Band and descendants of Frances Johnson as partners. The community should be contacted first, and their wishes on access and use respected.
\end{document}